% D9-1.0
% Auto-generated from markdown source
% origin: ai

\documentclass[12pt, a4paper]{article}
\usepackage[utf8]{inputenc}
\usepackage[T1]{fontenc}
\usepackage{amsmath, amssymb, amsfonts, amsthm}
\usepackage{textcomp}
\usepackage{geometry}
\usepackage{hyperref}
\usepackage{booktabs}
\usepackage{graphicx}
\usepackage{fancyhdr}
\usepackage{tcolorbox}
\usepackage{enumitem}

\geometry{margin=1in}
\pagestyle{fancy}
\fancyhf{}
\rhead{D9-1.0}
\lhead{\thepage}

\newtheorem{definition}{Definition}[section]
\newtheorem{theorem}{Theorem}[section]
\newtheorem{corollary}{Corollary}[section]
\newtheorem{proposition}{Proposition}[section]
\newtheorem{lemma}{Lemma}[section]
\newtheorem{conjecture}{Conjecture}[section]

\newcommand{\calI}{\mathcal{I}}
\newcommand{\calP}{\mathcal{P}}
\newcommand{\calE}{\mathcal{E}}
\newcommand{\calS}{\mathcal{S}}
\newcommand{\calH}{\mathcal{H}}
\newcommand{\calO}{\mathcal{O}}
\newcommand{\calF}{\mathcal{F}}
\newcommand{\calU}{\mathcal{U}}
\newcommand{\calV}{\mathcal{V}}
\newcommand{\calA}{\mathcal{A}}
\newcommand{\calG}{\mathcal{G}}
\newcommand{\calB}{\mathcal{B}}
\newcommand{\calC}{\mathcal{C}}
\newcommand{\calL}{\mathcal{L}}
\newcommand{\calM}{\mathcal{M}}
\newcommand{\calX}{\mathcal{X}}
\newcommand{\calY}{\mathcal{Y}}
\newcommand{\calW}{\mathcal{W}}
\newcommand{\calT}{\mathcal{T}}
\newcommand{\calZ}{\mathcal{Z}}
\newcommand{\calR}{\mathcal{R}}
\newcommand{\calQ}{\mathcal{Q}}
\newcommand{\calD}{\mathcal{D}}
\newcommand{\calN}{\mathcal{N}}
\newcommand{\calK}{\mathcal{K}}

\newcommand{\Tr}{\text{Tr}}
\newcommand{\E}{\mathbb{E}}
\newcommand{\diff}{\mathrm{d}}
\newcommand{\KL}{D_{\mathrm{KL}}}
\newcommand{\argmax}{\arg\max}
\newcommand{\argmin}{\arg\min}
\newcommand{\hbarF}{\hbar}

\newtcolorbox{thmbox}[1]{
  colback=blue!5!white, colframe=blue!75!black, fonttitle=\bfseries,
  title={#1}, breakable
}
\newtcolorbox{defbox}[1]{
  colback=green!5!white, colframe=green!75!black, fonttitle=\bfseries,
  title={#1}, breakable
}

\begin{document}

\title{D9-1.0}
\author{RALPH Loop}
\date{}
\maketitle

\begin{abstract}
This paper presents D9-1.0, formalizing the underlying principles using information-theoretic methods.

\textbf{Keywords}: D9-1.0, information theory, formalization.
\end{abstract}

\tableofcontents
\section{D9-1.0}


****D9-1.0-r01extended-framework-integration
****D9  + C1-D9-IV + D5-Chimera + D2/D9/D1-S12 + D1-S12


\#\#

****D9
\[\text{Intelligence} \equiv \mathcal{C}_1 \land \mathcal{C}_2 \land \mathcal{C}_3 \land \mathcal{C}_4 \land \mathcal{C}_5\]

\textbf{*}***


\#\#

\subsubsection{ 1}


D9 ****

|  |  |  |
|------|----------|---------------|
| \(\mathcal{C}_1\) | OOD  | C1- |
| \(\mathcal{C}_2\) |  | D5 |
| \(\mathcal{C}_3\) |  | D1-S12 chimera-N  |
| \(\mathcal{C}_4\) |  | C1-D9-IV- |
| \(\mathcal{C}_5\) |  | D5-Chimera- |


\subsubsection{ 2}


\textbf{ D2/D9/D1-S12 }

\textbf{ D9-1.0-C6-}
\[\mathcal{C}_6 \equiv |A| > I_{plan}^* + \delta_{margin} \land t > t_c\]

\textbf{*}*** +
**** \(\Psi_{\text{intelligence}} = \prod_{i=1}^{4} \Phi\left(\frac{X_i - X_i^c}{X_i^c}\right)\)


\subsubsection{ 3}


\textbf{ D1-S12  + D5-Chimera }

\textbf{ D9-1.0-C7}
\[\mathcal{C}_7 \equiv N \geq N_c \land \mathcal{H} > \mathcal{H}_c\]

****
\begin{itemize}
\item \(N_c \approx 500\)D1-S12-T1
\item \(N_{quantum} \approx 200-300\)D5-C-T4
\item \(\mathcal{H}\)  chimera D5-C-T2
\end{itemize}


\textbf{*}***


\#\#

\#\#\#

\textbf{D9-1.0 }
\[\text{Intelligence}(A) \equiv \bigwedge_{i=1}^{5} \mathcal{C}_i \land \mathcal{C}_6 \land \mathcal{C}_7\]

****
\[\text{Intelligence} = \mathcal{C}_{\text{basic}} \land \mathcal{C}_{\text{capacity-time}} \land \mathcal{C}_{\text{cognitive}}\]


\#\#

\#\#\#

\textbf{D2/D9/D1-S12}

\textbf{ sketch}
1. \textbf{\(\mathcal{C}_6\)  \(\mathcal{C}_1\)-C5}\(\mathcal{C}_6\)
2. \textbf{\(\mathcal{C}_7\)  \(\mathcal{C}_4\)} \(\mathcal{T}\)  \(N < N_c\)  gradient\_locking
3. \textbf{\(\mathcal{C}_6 \land \mathcal{C}_7\) }- N

\textbf{**D9-1.0 }**


\subsection{ S18}


\#\#\#

****
\[\frac{d}{dt} Q[\psi] = 0, \quad Q = \text{} \equiv \text{}\]

\textbf{D9-1.0 }
\[\frac{\partial \psi}{\partial t} = -\nabla_\psi \left[ \mathcal{L}_{\text{basic}} + \mathcal{L}_{\text{capacity}} + \mathcal{L}_{\text{cognitive}} \right]\]

****
\begin{itemize}
\item \(\mathcal{L}_{\text{basic}}\) +  +  +  +
\item \(\mathcal{L}_{\text{capacity}} = \alpha \cdot \max(0, I_{plan}^* + \delta - |A|) + \eta \cdot \max(0, t_c - t)\)
\item \(\mathcal{L}_{\text{cognitive}} = \gamma \cdot \max(0, N_c - N) + \delta \cdot \max(0, \mathcal{H}_c - \mathcal{H})\)
\end{itemize}


****
\[Q_{\text{total}} = Q_{\text{gen}} + Q_{\text{cap}} + Q_{\text{cog}}\]


\#\#

\#\#\#

```
D9-1.0


1.
1.1 D9
1.2
1.3

2. C1-C5
2.1-2.5

3. C6
3.1 -
3.2
3.3 D2-1/D2-2

4. C7
4.1 N500 D1-S12-T1
4.2 Chimera D5-C-T2
4.3 -D5-C-T4

5.
5.1
5.2

6. S18
6.1
6.2  Q\_total

7.
7.1 C1-D9-IV
7.2 D5-Chimera -
7.3 D1-S12 N

8.
```


\#\#

1. \textbf{ D9 }****
2. ****C1-D9-IV + D5-Chimera + D2/D9/D1-S12
3. ****
4. ****


\#\#

\begin{itemize}
\item [x]
\item [x]  S18
\item [x]
\item [ ]
\end{itemize}



* D9-1.0 *
*2026-05-04 | r01 | D9 *

\begin{thebibliography}{99}
\bibitem{ref1} Shannon, C.E. (1948). A Mathematical Theory of Communication.
\bibitem{ref2} Cover, T.M. (2006). Elements of Information Theory.
\end{thebibliography}

\end{document}